% !PRINTED 194
\weberpage{194}
\weq{8}{
\begin{array}{ll}
\left(y\sqrt{\dfrac{\varkappa}{\varkappa'}}\right)^{n^2-4}\! A\!\left(\dfrac{z^2}{x^2y^2}\right) = A(x^2), & A\!\left(\dfrac{1}{x^2}\right) = n\sqrt{\dfrac{\varkappa'}{\varkappa}}^{\,n^2-4}, \\[8pt]
\left(y\sqrt{\dfrac{\varkappa}{\varkappa'}}\right)^{n^2}\! B\!\left(\dfrac{z^2}{x^2y^2}\right) = (-1)^{\frac{n}{2}} B(x^2), & B\!\left(\dfrac{1}{x^2}\right) = (-1)^{\frac{n}{2}}\sqrt{\dfrac{\varkappa'}{\varkappa}}^{\,n^2-4}, \\[8pt]
\left(y\sqrt{\dfrac{\varkappa}{\varkappa'}}\right)^{n^2}\! C\!\left(\dfrac{z^2}{x^2y^2}\right) = C(x^2), & C\!\left(\dfrac{1}{x^2}\right) = \sqrt{\dfrac{\varkappa'}{\varkappa}}^{\,n^2-4}, \\[8pt]
\left(y\sqrt{\dfrac{\varkappa}{\varkappa'}}\right)^{2}\! D\!\left(\dfrac{z^2}{x^2y^2}\right) = (-1)^{\frac{n}{2}} D(x^2), & D\!\left(\dfrac{1}{x^2}\right) = (-1)^{\frac{n}{2}}\sqrt{\dfrac{\varkappa'}{\varkappa}}^{\,n^2-4}, \\[8pt]
& n \text{ gerade.} % [sic] exponent printed as 2, not n^2, on the D-row of this system
\end{array}
}
\weq{9}{
\begin{array}{ll}
\left(y\sqrt{\dfrac{\varkappa}{\varkappa'}}\right)^{n^2-1}\! A\!\left(\dfrac{z^2}{x^2y^2}\right) = (-1)^{\frac{n-1}{2}} C(x^2), & A\!\left(\dfrac{1}{x^2}\right) = (-1)^{\frac{n-1}{2}}\sqrt{\dfrac{\varkappa'}{\varkappa}}^{\,n^2-1}, \\[8pt]
\left(y\sqrt{\dfrac{\varkappa}{\varkappa'}}\right)^{n^2-1}\! B\!\left(\dfrac{z^2}{x^2y^2}\right) = D(x^2), & B\!\left(\dfrac{1}{x^2}\right) = \sqrt{\dfrac{\varkappa'}{\varkappa}}^{\,n^2-1}, \\[8pt]
\left(y\sqrt{\dfrac{\varkappa}{\varkappa'}}\right)^{n^2-1}\! C\!\left(\dfrac{z^2}{x^2y^2}\right) = (-1)^{\frac{n-1}{2}} A(x^2), & C\!\left(\dfrac{1}{x^2}\right) = (-1)^{\frac{n-1}{2}}\, n\, \sqrt{\dfrac{\varkappa'}{\varkappa}}^{\,n^2-1}, \\[8pt]
\left(y\sqrt{\dfrac{\varkappa}{\varkappa'}}\right)^{n^2-1}\! D\!\left(\dfrac{z^2}{x^2y^2}\right) = B(x^2), & D\!\left(\dfrac{1}{x^2}\right) = \sqrt{\dfrac{\varkappa'}{\varkappa}}^{\,n^2-1}, \\[8pt]
& n \text{ ungerade.}
\end{array}
}

Nach diesen Formeln kann man für den Fall eines ungeraden $n$ die vier
Polynome $A(x^2)$, $B(x^2)$, $C(x^2)$, $D(x^2)$ auf eines von ihnen, z. B. auf
$A(x^2)$, zurückführen.

Wenn man von den Formeln I, II je die drei ersten durch die letzte
dividiert, so erhält man die Multiplikationsformeln für die elliptischen
Funktionen (\S\,42).

\begin{center}
III. Bei geradem $n$: \qquad\qquad IV. Bei ungeradem $n$:
\end{center}
\[
\begin{array}{ll}
\mathrm{sn}\,nv = \dfrac{xyzA(x^2)}{D(x^2)}, & \mathrm{sn}\,nv = \dfrac{xA(x^2)}{D(x^2)}, \\[10pt]
\mathrm{cn}\,nv = \dfrac{B(x^2)}{D(x^2)}, & \mathrm{cn}\,nv = \dfrac{yB(x^2)}{D(x^2)}, \\[10pt]
\mathrm{dn}\,nv = \dfrac{C(x^2)}{D(x^2)}, & \mathrm{dn}\,nv = \dfrac{zC(x^2)}{D(x^2)}.
\end{array}
\]

Die Koeffizienten von $A$, $B$, $C$, $D$ hängen noch von $\omega$ oder von
$\varkappa^2$ ab. Über die Art dieser Abhängigkeit geben die Rekursionsformeln
Aufschluß, die man zur sukzessiven Berechnung von
%CONT>>
